% Gerado a partir de MONOGRAFIA_PASE.md.
\begin{apendicesenv}
\partapendices

\chapter[Prompts versionados da prova de conceito]{\textit{Prompts} versionados da prova de conceito}
A prova de conceito delega a modelos de linguagem apenas decisões específicas, cada uma definida por um \textit{prompt} próprio, mantido como arquivo versionado no repositório, e nunca como texto embutido no código. O Quadro~\ref{quadro:A1} relaciona os quinze \textit{prompts} à camada e ao componente da Arquitetura PASE que cada um implementa.

\begin{quadro}[htb]
\centering
\caption{Prompts da prova de conceito e componentes correspondentes}
\label{quadro:A1}
\footnotesize
\renewcommand{\arraystretch}{1.25}
\begin{tabularx}{\textwidth}{|X|X|X|}
\hline
\textbf{\textit{Prompt}} & \textbf{Função} & \textbf{Camada e componente} \\
\hline
P0 & Extração de intenção a partir de tickets de suporte & Tarefas: captação passiva \\
\hline
P0a & Agrupamento de tickets por tema & Tarefas: captação passiva \\
\hline
P1 & Normalização da demanda & Tarefas: \textit{middleware} de normalização \\
\hline
P1b & Avaliação de Zona Dinâmica & Tarefas: Avaliador de Zonas Dinâmicas \\
\hline
P2a & Coerência sistêmica com o contexto do produto & Tarefas: Motor de Filtragem \\
\hline
P2b & Qualificação por impacto, complexidade, confiança e risco & Tarefas: Motor de Filtragem \\
\hline
P3 & Geração exploratória de propostas & Tarefas: fluxo da inteligência artificial \\
\hline
P3b & Análise de lacunas entre módulos & Tarefas: fluxo da inteligência artificial \\
\hline
P4 & Planejamento técnico e critérios de aceite & Execução: fase de planejamento \\
\hline
P5 & Geração e correção de código & Execução: fase de execução \\
\hline
P6 & Geração de \textit{payload} de Zona Dinâmica & Execução: UPA \\
\hline
P9 & Análise de deriva de contexto dos descartes & Tarefas: Reavaliação Cíclica \\
\hline
P10 & Varredura de auditoria do código & Orquestração: Ciclo de Auditoria \\
\hline
P11 & Proposta de atualização de um documento de contexto & Orquestração: homologação \\
\hline
P12 & Revisão independente da entrega & Execução: revisão independente \\
\hline
\end{tabularx}
\fonte{elaborado pelos autores (2026).}
\end{quadro}

\chapter[Contrato canônico das Zonas Dinâmicas]{Contrato canônico das Zonas Dinâmicas}
O contrato rígido (\textit{Hard Constraint}) que as Zonas Dinâmicas aplicam é um único documento canônico em JSON Schema, validado no servidor antes de persistir cada \textit{payload} e espelhado no cliente, onde ocorre a validação decisiva. Um teste automatizado compara os dois campo a campo. O contrato admite apenas três componentes nativos, rótulos de até quarenta caracteres, no máximo três itens e ligações restritas às rotas internas do produto. Qualquer propriedade não prevista invalida o \textit{payload}.

{\scriptsize
\begin{verbatim}
{
  "$schema": "https://json-schema.org/draft/2020-12/schema",
  "title": "Contrato das Zonas Dinâmicas",
  "type": "object",
  "additionalProperties": false,
  "required": ["widget", "label"],
  "properties": {
    "widget": { "type": "string", "enum": ["shortcut_list", "summary_card", "alert_banner"] },
    "label": { "type": "string", "minLength": 1, "maxLength": 40 },
    "maxItems": { "type": "integer", "minimum": 1, "maximum": 3 },
    "tone": { "type": "string", "enum": ["neutral", "atencao"] },
    "items": {
      "type": "array",
      "maxItems": 3,
      "items": {
        "type": "object",
        "additionalProperties": false,
        "required": ["label", "href"],
        "properties": {
          "label": { "type": "string", "minLength": 1, "maxLength": 40 },
          "href": {
            "type": "string",
            "pattern": "^/(dashboard|transacoes|relatorios|orcamentos|metas|suporte)$"
          }
        }
      }
    }
  },
  "allOf": [
    {
      "if": { "properties": { "widget": { "enum": ["shortcut_list", "summary_card"] } } },
      "then": { "required": ["maxItems", "items"] }
    },
    {
      "if": { "properties": { "widget": { "const": "alert_banner" } } },
      "then": { "required": ["tone"] }
    }
  ]
}
\end{verbatim}
}

\chapter[Reprodução da prova de conceito]{Reprodução da prova de conceito}
A prova de conceito é composta por cinco serviços executados localmente em contêineres: o banco de dados do PASE, o núcleo do PASE, o painel de controle e os dois serviços do FinTrack. O repositório reúne o código dos dois sistemas, os roteiros de demonstração e as suítes de testes. [Indicar aqui como a banca pode acessar o repositório, que atualmente é privado.]

A partir de um ambiente com Docker e com as chaves de acesso ao provedor de modelos configuradas em um arquivo de ambiente local, a sequência de reprodução é a seguinte. O comando \texttt{make reset} recria o ambiente do zero e popula o FinTrack e a base inicial do PASE. Os comandos \texttt{make test-core} e \texttt{make test-client} executam as suítes do núcleo e da interface de programação do FinTrack, e a suíte da interface de usuário é executada no contêiner correspondente. O comando \texttt{make demos} executa os oito roteiros de demonstração do Quadro~\ref{quadro:4}, e \texttt{make metrics} exporta os arquivos dos quais saem as métricas da Tabela~\ref{tab:2}. Como os roteiros fazem chamadas reais a modelos de linguagem, sua execução tem custo e pode produzir resultados ligeiramente diferentes a cada vez.

\end{apendicesenv}
